\section{Modulación pasabanda $M$-aria y su espectro de potencia}

Para introducir estas modulaciones es conveniente recurrir a la nocion
de \emph{envolvente compleja} que fue ampliamente utilizada
para la modulación pasabanda analógica. Recordamos que se plantea
la relación
\[x_c(t) = \Re\left[x_{ce}(t)\e^{j\omega_ct}\right], \]
donde $x_{ce}(t)=x_i(t)+jx_q(t) $ es una se\~nal compleja de banda base que juega el rol de un
fasor, variando en el tiempo lentamente respecto de las variaciones de
la portadora.

En el caso de la modulación digital, se considera una envolvente compleja
\[
x_{ce}(t) =  \sum_k  \underbrace{(a_k + j b_k)}_{c_k} p(t-kT_s)
% =  \underbrace{\sum a_k p(t-kT_s)}_{x_i(t)} + j \underbrace{\sum b_k p(t-kT_s)}_{x_q(t)}.
\]
La expresión anterior
es una generalización de la onda PAM a señales complejas: un tren de pulsos
multiplicados por símbolos complejos $c_k=a_k + j b_k$, que
varían en un conjunto discreto llamado
 {\em constelación de símbolos }.
%
\begin{minipage}{0.45\textwidth}
    \bpgf{Imagenes/MDPasa/MDPasa-3}
    \begin{tikzpicture}
        \draw (-3,0) -- (3,0) node [anchor=north] {$i$};
        \draw (0,-3) -- (0,3) node [anchor=west] {$q$};
        \draw [fill] (2,1) node [anchor=west] {$c_k$} circle (0.05);
        \draw [fill] (1,2) circle (0.08);
        \draw [fill] (-1,2) circle (0.08);
        \draw [fill] (-2,1) circle (0.08);
        \draw [fill] (-2,-1) circle (0.08);
        \draw [fill] (-1,-2) circle (0.08);
        \draw [fill] (1,-2) circle (0.08);
        \draw [fill] (2,-1) circle (0.08);
    \end{tikzpicture}
    \endpgfgraphicnamed
\end{minipage}%\hfill
%\begin{minipage}{0.45\textwidth}
%    Constelación de símbolos.
%\end{minipage}

La onda pasabanda correspondiente a esta envolvente compleja es

\[x_c(t) = \Re\left[x_{ce}(t)\e^{j\omega_ct}\right] = \sum_{k}p(t-kT_s)[a_k \cos(\omega_c t) - b_k \sen(\omega_c t)].
\]


\subsection*{Ejemplos de constelaciones}
       \FloatBarrier
\begin{figure}[htb]\centering\begin{minipage}{\textwidth}
        \bpgf{Imagenes/MDPasa/MDPasa-4}
        \begin{tikzpicture}
            \begin{scope}[shift={(-3,0)}]
                \draw (-2,0) -- (2,0);
                \draw (0,-2) -- (0,2);
                \draw [fill] (0,0) node [anchor=north east] {$0$} circle (0.08);
                \draw [fill] (1,0) node [anchor=north] {$A$} circle (0.08);
                \path (1,1) node {ASK};
                \begin{scope}[shift={(0,-5)}]
                    \draw (-2,0) -- (2,0);
                    \draw (0,-2) -- (0,2);
                    \draw [fill] (0,1) circle (0.08);
                    \draw [fill] (0,-1) circle (0.08);
                    \draw [fill]  (1,0) circle (0.08);
                    \draw [fill] (-1,0) circle (0.08);
                    \path (1,1) node {4-PSK};
                \end{scope}
                \begin{scope}[shift={(0,-10)}]
                    \draw (-2,0) -- (2,0);
                    \draw (0,-2) -- (0,2);
                    \draw [fill] (1,1) node [anchor=west] {$A+jA$} circle (0.08);
                    \draw [fill] (1,-1) circle (0.08);
                    \draw [fill]  (-1,1) circle (0.08);
                    \draw [fill] (-1,-1) circle (0.08);
                    \path (-2,2) node {4QAM};
                \end{scope}
            \end{scope}
            \begin{scope}[shift={(3,0)}]
                \draw (-2,0) -- (2,0);
                \draw (0,-2) -- (0,2);
                \draw [fill] (-1,0) node [anchor=north] {$-A$} circle (0.08);
                \draw [fill] (1,0) node [anchor=north] {$A$} circle (0.08);
                \path (1,1) node {BPSK};
                \begin{scope}[shift={(0,-5)}]
                    \draw (-2,0) -- (2,0);
                    \draw (0,-2) -- (0,2);
                    \foreach \x in {0,45,...,315}
                        \draw [fill] (\x:1) circle (0.08);
                    \draw (0,0) circle (1);
                    \path (2,1) node {8-PSK};
                \end{scope}
                \begin{scope}[shift={(0,-10)}]
                    \draw (-2,0) -- (2,0);
                    \draw (0,-2) -- (0,2);
                    \foreach \x in {-1.5,-0.5,...,1.5}
                        \foreach \y in {-1.5,-0.5,...,1.5}
                            \draw [fill] (\x,\y) circle (0.08);
                    \path (-2,2) node {16QAM};
                \end{scope}
            \end{scope}
        \end{tikzpicture}
        \endpgfgraphicnamed \end{minipage}
\end{figure}
%\FloatBarrier

\subsection*{Densidad espectral de potencia, onda pasabanda $M$-aria}

En el caso general $M$-ario, se puede obtener también la densidad espectral de potencia  a partir de
la de su envolvente compleja, generalizando la fórmula (\ref{eq:dep}).
\[
    \col{S_{x_c}(f) = \frac{1}{4}S_{x_{ce}}(f-f_c) +
    \frac{1}{4}S_{x_{ce}}(f+f_c)}
\]
La pregunta es entonces como calcular la densidad espectral de potencia para la onda \emph{compleja} de
banda base
\[x_{ce}(t) = \sum \underbrace{(a_k+j b_k)}_{c_k} p(t-kT_s),\]
donde suponemos que los símbolos $\{c_k\}$ son variables aleatorias complejas, independientes e
identicamente distribuidas. Suponemos, sin perder generalidad, que $E[c_k]=0$, es decir que la
costelación es balanceada alrededor del cero, con puntos equiprobables.

\newpage

La respuesta es que se puede utilizar la misma fórmula de antes (con $\mu=0$):
\[
    S_{x_{ce}} = \sigma_c^2 f_s |\hat{P}(f)|^2,
        \]
donde ahora $\sigma_c^2 = E[|c_k|^2] = E[a_k^2 + b_k^2]$, varianza del módulo de los símbolos.
Con este método consideramos los dos casos principales.

\subsubsection*{Modulación MPSK}


En este caso la constelación se define en coordenadas polares:
\[
c_k \in A e^{j(\theta_0 + \frac{2\pi}{M}{m})} \quad m=0,1,\ldots,M-1.\]

%\begin{minipage}
\FloatBarrier
\begin{figure}[htb]
    \centering\begin{minipage}{\textwidth}
    \bpgf{Imagenes/MDPasa/MDPasa-8}
    \begin{tikzpicture}
        \begin{scope}[shift={(5,0)}]
            \draw (-1.5,0) -- (1.5,0);
            \draw (0,-1.5) -- (0,1.5);
            \draw [fill] (1,0) circle (0.07);
            \draw [fill] (-1,0) circle (0.07);
            \draw [fill] (0,1) circle (0.07);
            \draw [fill] (0,-1) circle (0.07);
            \path (0.5,1) node [anchor=west] {\small $\theta_0=0, M=4$};
        \end{scope}
        \begin{scope}[shift={(5,-4)}]
            \draw (-1.5,0) -- (1.5,0);
            \draw (0,-1.5) -- (0,1.5);
            \draw [fill] (1,1) circle (0.07);
            \draw [fill] (-1,1) circle (0.07);
            \draw [fill] (1,-1) circle (0.07);
            \draw [fill] (-1,-1) circle (0.07);
            \path (1.5,1) node [anchor=west] {\small $\theta_0=\frac{\pi}{4}, M=4$};
            \path (1.5,0.5) node [anchor=west] {\footnotesize (equivale a 4 QAM)};
        \end{scope}
            \begin{scope}[shift={(5,-8)}]
            \draw (-1.5,0) -- (1.5,0);
            \draw (0,-1.5) -- (0,1.5);
            \foreach \x in {0,45,...,315}
                \draw [fill] (\x:1) circle (0.07);
            \draw (0,0) -- node [above=0.1mm] {\footnotesize $A$} (45:1);
            \path (1,1) node [anchor=west] {\small $\theta_0=0, M=8$};
        \end{scope}
    \end{tikzpicture}
    \endpgfgraphicnamed\end{minipage}
    \end{figure}
%\FloatBarrier
%\end{minipage}

\noindent Suponiendo que los puntos son equiprobables, tenemos $\mu=0$ y $\sigma_c^2 = E[|c_k|^2] = A^2$.
Por tanto
\[\begin{gathered}
            S_{x_{ce}}(f) = A^2 f_s |\hat{P}(f)|^2
            \Rightarrow S_{x_c}(f) = \frac{A^2 f_s}{4}\left(|\hat{P}(f-f_c)|^2+|\hat{P}(f+f_c)|^2\right),
        \end{gathered}\]% \right\rbrace\,\text{igual que BPSK.}
que es el mismo resultado que en BPSK, independiente del número de puntos $M$ de la constelación.

\newpage

\subsubsection*{Modulación QAM}

Aquí la constelación se describe en coordenadas cartesianas, $c_k = a_k+j b_k$, donde
cada coordenada toma $M_a$, $M_b$ valores polares, $M=M_a M_b$.

\begin{figure}[htb]
    \centering\begin{minipage}{\textwidth}
    \bpgf{Imagenes/MDPasa/MDPasa-9}
    \begin{tikzpicture}
        \draw [->] (-4,0) -- (4,0);
        \draw [->] (0,-4) -- (0,4);
        \foreach \x in {-3,-1,...,3}
            {\draw [dotted] (\x,-3) -- (\x,3);
            \draw [dotted] (-3,\x) -- (3,\x);
            \foreach \y in {-3,-1,...,3}
                \draw [fill] (\x,\y) circle (0.07);}
        \path (1,0) node [anchor=north west] {$A$} -- (3,0) node [anchor=north west] {$3A$};
    \end{tikzpicture}
    \endpgfgraphicnamed\end{minipage}
\end{figure}
\FloatBarrier


Tenemos entonces que $\sigma_c^2 = \sigma_a^2+\sigma_b^2$, y cada uno de estos términos se calcula como
\[
\sigma_a^2 = \frac{2}{M_a}\left(A^2 +
(3A)^2+\cdots+(M_a-1)^2A^2\right) = \frac{M_a^2-1}{3}A^2.\]
Tenemos:
\begin{align*}
S_{x_{ce}}(f) &= (\sigma_a^2+\sigma_b^2) f_s |\hat{P}(f)|^2 \\
\Rightarrow S_{x_c}(f) &=
\frac{(\sigma_a^2+\sigma_b^2)f_s\left(|\hat{P}(f-f_c)|^2 +
|\hat{P}(f+f_c)|^2\right)}{4}.
\end{align*}



Habitualmente $M_a = M_b = \sqrt{M}$. En ese caso
$\sigma_a^2+\sigma_b^2 = \frac{2(M-1)}{3}A^2$.





